\documentclass[11pt]{article}
\usepackage[margin=1in]{geometry}
\usepackage{amsmath,amssymb,amsthm,hyperref}
\usepackage{xurl}
\newtheorem{lemma}{Lemma}
\newcommand{\F}{\mathbb F}
\newcommand{\rank}{\operatorname{rank}}
\newcommand{\row}{\operatorname{row}}
\title{Independent Audit of the Affine-Slice Row-Prefix Character Formula}
\author{Astra P3 continuation: independent algebraic audit}
\date{28 October 2026}
\begin{document}\maketitle

\paragraph{Audited claim.}
The new \path|affine_sampler/affine_slice_prefix_sampler.py| counts prefixes
of acquired right-matrix rows at cost $2^c$ times a polynomial, where $c$ is
the original binary affine-slice codimension. Partially fixing a row does
not add an exponential Fourier factor. The following derivation independently
supports the implemented three-mass formula. It concerns acquired-coordinate
row prefixes, not arbitrary original-input pins.

\section{Complete rows and quotient factors}

Let $F=\F_q$, $q=2^a$, $X\in F^{p\times d}$, $Y\in F^{d\times r}$,
and $XY=Z$. For a frequency matrix $D\in F^{d\times r}$ use
\[
 \chi_D(Y)=(-1)^{\operatorname{Tr}_{F/\F_2}
                         \sum_i D_iY_i^{\mathsf T}}.
\]
Fix the first $t$ rows of $Y$ to a matrix $A$ and put
$R=\row A$, $\alpha=\dim R$. Write $X=[X_0\ X_1]$ and
$Y=[A;B]$. Passing to $F^r/R$ gives
\[
 X_1\overline B=\overline Z.
\]
For every solution of that quotient equation, $X_0A=Z-X_1B$ has
$q^{p(t-\alpha)}$ solutions. Each quotient row of $B$ has $q^\alpha$
lifts. Summing its character over those lifts is zero unless the
corresponding row of $D$ lies in $R^\perp$, and otherwise gives $q^\alpha$.
Therefore, when all remaining frequency rows annihilate $R$, the signed
sum is
\begin{equation}\label{eq:quotient}
 \chi_{D_{<t}}(A)\,
 q^{\alpha(d-t)+p(t-\alpha)}
 F_{p,d-t,r-\alpha,q}(k,\ell,\rho),
\end{equation}
where $F$ is the inherited signed universal product-fiber formula and
\[
 k=\rank[A;Z]-\alpha,\quad
 \ell=\rank D_{\ge t},\quad
 \rho=\rank(ZD_{\ge t}^{\mathsf T}).
\]
The latter two ranks may be computed in the original ambient coordinates:
the frequency rows annihilate $R$, so they represent quotient functionals
faithfully. If $r=\alpha$, the quotient has dimension zero and the
corresponding scalar is $F_{p,d-t,0,q}(0,0,0)=q^{p(d-t)}$; the same formula
holds. All powers in the prefactor are nonnegative since $t\ge\alpha$.

\section{A partially fixed next row}

Suppose the first $s$ binary coordinates of the next row are fixed to
$b\in\F_2^s$. The Fourier indicator is
\[
 2^{-s}\sum_{\lambda\in\F_2^s}
       (-1)^{\lambda\cdot b}\,
       (-1)^{\operatorname{Tr}(\delta_\lambda Y_t^{\mathsf T})},
\]
where $\delta_\lambda$ uses the actual trace-dual basis of $F$. It changes
the first remaining frequency to $w=D_t+\delta_\lambda$. No assumption that
the acquired field identity has coordinate word one is valid or used.

Let $D_{>t}$ contain the later frequency rows. If one fails to annihilate
$R$, the sum is zero independently of $\lambda$. Otherwise define
\[
 K=\row D_{>t},\quad E=R^\perp,\quad
 L=E\cap\{w:Zw^{\mathsf T}\in ZK^{\mathsf T}\}.
\]
Then $K\subseteq L\subseteq E$. Put
$\ell_0=\dim K$ and $\rho_0=\dim ZK^{\mathsf T}$.
The only possible rank pairs for the remaining frequency matrix are
\[
\begin{array}{c|c}
 w\text{ lies in}&(\ell,\rho)\\\hline
 K&(\ell_0,\rho_0)\\
 L\setminus K&(\ell_0+1,\rho_0)\\
 E\setminus L&(\ell_0+1,\rho_0+1).
\end{array}
\]
Values outside $E$ vanish in the lift sum.

For $T\in\{K,L,E\}$ define the signed affine mass
\[
 m_T=\sum_{\lambda:D_t+\delta_\lambda\in T}(-1)^{\lambda\cdot b}.
\]
Membership in $T$ is a binary linear system in $\lambda$. If its solution
set is empty, $m_T=0$. If it is $\lambda_0+V$ and the character is
nontrivial on $V$, again $m_T=0$. Otherwise
$m_T=(-1)^{\lambda_0\cdot b}2^{\dim V}$. Thus each of the three masses
is computed by one binary elimination, without enumerating $2^s$ characters.

Combining this partition with \eqref{eq:quotient} gives exactly
\begin{align*}
 &2^{-s}\chi_{D_{<t}}(A)
       q^{\alpha(d-t)+p(t-\alpha)}\\
 &\quad\times\bigl[
 m_KF(k,\ell_0,\rho_0)
 +(m_L-m_K)F(k,\ell_0+1,\rho_0)
 +(m_E-m_L)F(k,\ell_0+1,\rho_0+1)\bigr],
\end{align*}
with the dimensions on $F$ those in \eqref{eq:quotient}. The expression
is an integer because it is the signed sum over the original partially
fixed matrix-fiber set. The implementation also checks this division.

The code acquires $K$ by the nullspace of $D_{>t}$, acquires $L$ by the
rows of $A$ together with $Z^{\mathsf T}$ applied to the annihilator of
$ZD_{>t}^{\mathsf T}$, and uses the rows of $A$ for $E$. These are exactly
the conditions above, including their transpose directions.

\section{Affine slice, exact sampling, and cost}

The original affine slice contributes its original $2^c$ independent
characters and their phases. Apply the preceding formula to each one and
divide by $2^c$. The number of characters does not increase with the prefix
length. At a completely fixed $Y$, the result is zero unless its affine
constraints hold and $\row Z\subseteq\row Y$; otherwise it is
$q^{p(d-\rank Y)}$. This agrees with solving the left linear fibers.

Successively choose right-matrix bits with probability equal to their exact
completion count divided by the parent count. Once $Y$ is complete, solve
each left-row affine system and draw uniformly from its kernel translate.
The prefix probabilities telescope: a right matrix is selected with its
number of left completions divided by the full population, and each
completion is then selected with the reciprocal of that number. Every
original matrix pair therefore has exactly equal probability. The acquired
coordinate maps and corrected original affine reconstruction are bijective,
so uniformity carries back to the original independent inputs.

There are at most $adr$ bit-prefix decisions. Each query uses $2^c$
characters, three polynomial-size binary solves per character, and the
polynomial signed-fiber scalar. Finite-field arithmetic is manufactured,
not an oracle; all integers have polynomial bit length. The bound is
$2^c\operatorname{poly}(a,p,d,r)$ per sample after coordinate acquisition.
The native gate frontend enforces $2^c$ no larger than the actual compact
input-description byte length, giving a polynomially admitted native family.
For unrestricted $c$, the explicit exponential dependence remains.
Random integer interval draws and inherited coordinate acquisition have
polynomial expectation; no universal worst-case random-bit bound is claimed.

\section{Independent executed check}

\path|prefix_character_independent_audit.py| enumerates actual matrix
factor pairs, computes product targets and characters directly, and compares
each one-character signed prefix sum with the implemented formula. It does
not use the affine character sum as its reference, so cancellation between
different original characters cannot conceal an error.

There were 14,496 exact signed prefix comparisons from 8,704 literal factor
pairs over four rectangular parameter cases. The binary $p=d=r=2$ case
checks every target and every original character. The other cases include
three-row prefixes, quotient dimension zero, partial field-coordinate rows,
rectangular left factors, and $\F_4$ represented in basis $(\omega,1)$,
whose identity coordinate is two. Every comparison passed. Exact data are
in \path|prefix_character_independent_audit_results.json|.

No mathematical or implementation defect was found in the audited core.
The result repairs the prior full-fiber rejection obstruction for arbitrary
affine slices at bounded/logarithmic acquired codimension. It does not by
itself remove the $2^c$ dependence or establish arbitrary original-input
conditional-count closure.
\end{document}
